\documentclass[10pt,a4paper]{amsart}
\usepackage{amsmath,amssymb,amsthm,graphicx}
% For dealing with references we use the comment environment
\usepackage{verbatim}
\newenvironment{reference}{\comment}{\endcomment}
%\newenvironment{reference}{}{}
\newenvironment{slogan}{\comment}{\endcomment}
\newenvironment{history}{\comment}{\endcomment}

% For commutative diagrams we use Xy-pic
\usepackage[all]{xy}
% Theorem environments.
%
\theoremstyle{plain}
\newtheorem{theorem}[subsection]{Theorem}
\newtheorem{proposition}[subsection]{Proposition}
\newtheorem{lemma}[subsection]{Lemma}

\theoremstyle{definition}
\newtheorem{definition}[subsection]{Definition}
\newtheorem{example}[subsection]{Example}
\newtheorem{exercise}[subsection]{Exercise}
\newtheorem{situation}[subsection]{Situation}

\theoremstyle{remark}
\newtheorem{remark}[subsection]{Remark}
\newtheorem{remarks}[subsection]{Remarks}

\numberwithin{equation}{subsection}

% Macros
%
\def\lim{\mathop{\mathrm{lim}}\nolimits}
\def\colim{\mathop{\mathrm{colim}}\nolimits}
\def\Spec{\mathop{\mathrm{Spec}}}
\def\Hom{\mathop{\mathrm{Hom}}\nolimits}
\def\Ext{\mathop{\mathrm{Ext}}\nolimits}
\def\SheafHom{\mathop{\mathcal{H}\!\mathit{om}}\nolimits}
\def\SheafExt{\mathop{\mathcal{E}\!\mathit{xt}}\nolimits}
\def\Sch{\mathit{Sch}}
\def\Mor{\mathop{\mathrm{Mor}}\nolimits}
\def\Ob{\mathop{\mathrm{Ob}}\nolimits}
\def\Sh{\mathop{\mathit{Sh}}\nolimits}
\def\NL{\mathop{N\!L}\nolimits}
\def\CH{\mathop{\mathrm{CH}}\nolimits}
\def\proetale{{pro\text{-}\acute{e}tale}}
\def\etale{{\acute{e}tale}}
\def\QCoh{\mathit{QCoh}}
\def\Ker{\mathop{\mathrm{Ker}}}
\def\Im{\mathop{\mathrm{Im}}}
\def\Coker{\mathop{\mathrm{Coker}}}
\def\Coim{\mathop{\mathrm{Coim}}}

% Boxtimes
%
\DeclareMathSymbol{\boxtimes}{\mathbin}{AMSa}{"02}

%
% Macros for moduli stacks/spaces
%
\def\QCohstack{\mathcal{QC}\!\mathit{oh}}
\def\Cohstack{\mathcal{C}\!\mathit{oh}}
\usepackage{cleveref}
\crefname{originallemma}{original lemma}{original lemmas}
\Crefname{originallemma}{Original lemma}{Original lemmas}
\makeatletter
\newlabel{lemma-mayer-vietoris-hom}{{08BW}{source}{}{sourcepointer-1}{}}
\newlabel{lemma-mayer-vietoris-hom@cref}{{[originallemma][0][]08BW}{[1][0][]source}{}{}{}}
\makeatother
\crefname{originallemma}{original lemma}{original lemmas}
\Crefname{originallemma}{Original lemma}{Original lemmas}
\makeatletter
\newlabel{lemma-section-RHom-over-U}{{08DK}{source}{}{sourcepointer-2}{}}
\newlabel{lemma-section-RHom-over-U@cref}{{[originallemma][0][]08DK}{[1][0][]source}{}{}{}}
\makeatother
\crefname{originallemma}{original lemma}{original lemmas}
\Crefname{originallemma}{Original lemma}{Original lemmas}
\makeatletter
\newlabel{lemma-unbounded-describe-higher-direct-images}{{0D5X}{source}{}{sourcepointer-3}{}}
\newlabel{lemma-unbounded-describe-higher-direct-images@cref}{{[originallemma][0][]0D5X}{[1][0][]source}{}{}{}}
\makeatother
\crefname{originallemma}{original lemma}{original lemmas}
\Crefname{originallemma}{Original lemma}{Original lemmas}
\makeatletter
\newlabel{topology-lemma-create-map-from-subcollection}{{0D5Q}{source}{}{sourcepointer-4}{}}
\newlabel{topology-lemma-create-map-from-subcollection@cref}{{[originallemma][0][]0D5Q}{[1][0][]source}{}{}{}}
\makeatother
\title{BBD gluing of unbounded derived objects over an arbitrary open cover with vanishing negative Ext}
\author{The Stacks Project Authors}
\date{stacks-project:cohomology.tex:local-snapshot-20261011; GNU FDL 1.2 or later}
\begin{document}
\maketitle
\noindent\textit{Editorial note.} The original theorem credits BBD Theorem3.2.4; this source proves the stated special case without boundedness. All source-local gluing lemmas used for the count are included as solution support, not separately counted.\par\medskip
\noindent\textit{Editorial note.} At11432 the invoked property is uniqueness of an already existing solution on the limit cover, established by lemma-uniqueness11187-11221; the printed finite-case reference is unnecessary for this uniqueness assertion.\par\medskip
\noindent\textit{Editorial note.} At11430-11435 the unique isomorphism is printed K\_alpha→C while the chosen representative has direction C→I\_alpha; the latter represents its inverse. The original directions are retained and explicitly identified.\par\medskip
\noindent\textit{Editorial note.} At11546 the final rho indices are reversed in print; the compatibility is the one explicitly imposed in11353-11356 and stated in the preceding composition. Source formulas are retained.\par\medskip
\noindent\textit{Editorial scope.} One original target; complete source-local proof chain included, standard derived-sheaf primitives explicitly delimited.\par\medskip
\noindent\textit{Allowed standard prerequisites (editorial).}\begin{itemize}
\item Standard derived sheaf operations. K-injective resolutions, extension by zero for open embeddings, compatibility of derived pushforward with restriction, and exactness of filtered colimits on stalks.
\item Derived Hom and Mayer-Vietoris. Derived internal Hom computes local Ext and H\textasciicircum{}0 computes derived-category Hom; two-open Mayer-Vietoris gives the exact Hom sequence.
\item Elementary topology of a subcollection of opens. A covering family stable under intersections up to unions is the inverse image of a basis under a continuous map; take the topology generated by this family on the underlying set.
\end{itemize}\medskip
\ifdefined\hypertarget\hypertarget{sourcepointer-1}{}\fi
\noindent\textit{Source reference.} Original lemma 08BW: focused-seeds/ACL/sources/stacks-snapshot-20261011/cohomology.tex, source lines 8166--8184. Standard two-open exact Hom sequence; the gluing object and its compatibility are constructed locally. Displayed pointer is the original stable Stacks tag.\par\medskip
\ifdefined\hypertarget\hypertarget{sourcepointer-2}{}\fi
\noindent\textit{Source reference.} Original lemma 08DK: focused-seeds/ACL/sources/stacks-snapshot-20261011/cohomology.tex, source lines 10580--10589. Standard identification of local derived Hom with H\textasciicircum{}0 of internal Hom. Displayed pointer is the original stable Stacks tag.\par\medskip
\ifdefined\hypertarget\hypertarget{sourcepointer-3}{}\fi
\noindent\textit{Source reference.} Original lemma 0D5X: focused-seeds/ACL/sources/stacks-snapshot-20261011/cohomology.tex, source lines 7941--7949. Standard cohomology-sheaf description for derived direct images. Displayed pointer is the original stable Stacks tag.\par\medskip
\ifdefined\hypertarget\hypertarget{sourcepointer-4}{}\fi
\noindent\textit{Source reference.} Original lemma 0D5Q: focused-seeds/ACL/sources/stacks-snapshot-20261011/topology.tex, source lines 376--390. Elementary topology change providing a basis, explicitly permitted and not a derived descent theorem. Displayed pointer is the original stable Stacks tag.\par\medskip
\begin{situation}
\label{situation-locally-given}
Let $(X, \mathcal{O}_X)$ be a ringed space. We are given
\begin{enumerate}
\item a collection of opens $\mathcal{B}$ of $X$,
\item for $U \in \mathcal{B}$ an object $K_U$ in $D(\mathcal{O}_U)$,
\item for $V \subset U$ with $V, U \in \mathcal{B}$ an isomorphism
$\rho^U_V : K_U|_V \to K_V$ in $D(\mathcal{O}_V)$,
\end{enumerate}
such that whenever we have $W \subset V \subset U$ with $U, V, W$ in
$\mathcal{B}$, then $\rho^U_W = \rho^V_W \circ \rho ^U_V|_W$.
\end{situation}

\noindent
We won't be able to prove anything about this without making more
assumptions. An interesting case is where $\mathcal{B}$ is a basis
for the topology on $X$. Another is the case where we have a morphism
$f : X \to Y$ of topological spaces and the elements of $\mathcal{B}$
are the inverse images of the elements of a basis for the topology of $Y$.

\medskip\noindent
In Situation \ref{situation-locally-given} a {\it solution}
will be a pair $(K, \rho_U)$ where $K$ is an object of $D(\mathcal{O}_X)$
and $\rho_U : K|_U \to K_U$, $U \in \mathcal{B}$
are isomorphisms such that
we have $\rho^U_V \circ \rho_U|_V = \rho_V$ for all $V \subset U$,
$U, V \in \mathcal{B}$. In certain cases solutions are unique.

\begin{theorem}[BBD gluing lemma]
\label{theorem-glueing-bbd-general}
\begin{reference}
Special case of \cite[Theorem 3.2.4]{BBD}
without boundedness assumption.
\end{reference}
In Situation \ref{situation-locally-given} assume
\begin{enumerate}
\item $X = \bigcup_{U \in \mathcal{B}} U$,
\item for $U, V \in \mathcal{B}$ we have
$U \cap V = \bigcup_{W \in \mathcal{B}, W \subset U \cap V} W$,
\item for any $U \in \mathcal{B}$ we have $\Ext^i(K_U, K_U) = 0$
for $i < 0$.
\end{enumerate}
Then there exists an object $K$ of $D(\mathcal{O}_X)$
and isomorphisms $\rho_U : K|_U \to K_U$ in $D(\mathcal{O}_U)$ for
$U \in \mathcal{B}$ such that $\rho^U_V \circ \rho_U|_V = \rho_V$
for all $V \subset U$ with $U, V \in \mathcal{B}$.
The pair $(K, \rho_U)$ is unique up to unique isomorphism.
\end{theorem}

\begin{lemma}
\label{lemma-glue}
Let $(X, \mathcal{O}_X)$ be a ringed space. Let $X = U \cup V$ be
the union of two open subspaces of $X$. Suppose given
\begin{enumerate}
\item an object $A$ of $D(\mathcal{O}_U)$,
\item an object $B$ of $D(\mathcal{O}_V)$, and
\item an isomorphism $c : A|_{U \cap V} \to B|_{U \cap V}$.
\end{enumerate}
Then there exists an object $F$ of $D(\mathcal{O}_X)$
and isomorphisms $f : F|_U \to A$, $g : F|_V \to B$ such
that $c = g|_{U \cap V} \circ f^{-1}|_{U \cap V}$.
Moreover, given
\begin{enumerate}
\item an object $E$ of $D(\mathcal{O}_X)$,
\item a morphism $a : A \to E|_U$ of $D(\mathcal{O}_U)$,
\item a morphism $b : B \to E|_V$ of $D(\mathcal{O}_V)$, 
\end{enumerate}
such that
$$
a|_{U \cap V}  = b|_{U \cap V} \circ c.
$$
Then there exists a morphism $F \to E$ in $D(\mathcal{O}_X)$
whose restriction to $U$ is $a \circ f$
and whose restriction to $V$ is $b \circ g$.
\end{lemma}

\begin{proof}
Denote $j_U$, $j_V$, $j_{U \cap V}$ the corresponding open immersions.
Choose a distinguished triangle
$$
F \to Rj_{U, *}A \oplus Rj_{V, *}B \to Rj_{U \cap V, *}(B|_{U \cap V})
\to F[1]
$$
where the map $Rj_{V, *}B \to Rj_{U \cap V, *}(B|_{U \cap V})$ is the
obvious one and where
$Rj_{U, *}A \to Rj_{U \cap V, *}(B|_{U \cap V})$
is the composition of
$Rj_{U, *}A \to Rj_{U \cap V, *}(A|_{U \cap V})$
with $Rj_{U \cap V, *}c$. Restricting to $U$ we obtain
$$
F|_U \to A \oplus (Rj_{V, *}B)|_U \to (Rj_{U \cap V, *}(B|_{U \cap V}))|_U
\to F|_U[1]
$$
Denote $j : U \cap V \to U$. Compatibility of restriction to opens and
cohomology shows that both
$(Rj_{V, *}B)|_U$ and $(Rj_{U \cap V, *}(B|_{U \cap V}))|_U$
are canonically isomorphic to $Rj_*(B|_{U \cap V})$.
Hence the second arrow of the last displayed diagram has
a section, and we conclude that the morphism $F|_U \to A$ is
an isomorphism. Similarly, the morphism $F|_V \to B$ is an
isomorphism. The existence of the morphism $F \to E$ follows
from the Mayer-Vietoris sequence for $\Hom$, see
Lemma \ref{lemma-mayer-vietoris-hom}.
\end{proof}

\begin{lemma}
\label{lemma-vanishing-and-glueing}
Let $f : (X, \mathcal{O}_X) \to (Y, \mathcal{O}_Y)$ be a morphism
of ringed spaces. Let $\mathcal{B}$ be a basis for the topology on $Y$.
\begin{enumerate}
\item Assume $K$ is in $D(\mathcal{O}_X)$ such that
for $V \in \mathcal{B}$ we have $H^i(f^{-1}(V), K) = 0$ for $i < 0$.
Then $Rf_*K$ has vanishing cohomology sheaves in negative degrees,
$H^i(f^{-1}(V), K) = 0$ for $i < 0$ for all opens $V \subset Y$, and
the rule $V \mapsto H^0(f^{-1}V, K)$ is a sheaf on $Y$.
\item Assume $K, L$ are in $D(\mathcal{O}_X)$ such that
for $V \in \mathcal{B}$ we have
$\Ext^i(K|_{f^{-1}V}, L|_{f^{-1}V}) = 0$ for $i < 0$.
Then $\Ext^i(K|_{f^{-1}V}, L|_{f^{-1}V}) = 0$ for $i < 0$
for all opens $V \subset Y$ and
the rule $V \mapsto \Hom(K|_{f^{-1}V}, L|_{f^{-1}V})$ is a sheaf on $Y$.
\end{enumerate}
\end{lemma}

\begin{proof}
Lemma \ref{lemma-unbounded-describe-higher-direct-images} tells us
$H^i(Rf_*K)$ is the sheaf associated to the presheaf
$V \mapsto H^i(f^{-1}(V), K) = H^i(V, Rf_*K)$.
The assumptions in (1) imply that $Rf_*K$ has vanishing cohomology
sheaves in degrees $< 0$. We conclude that for any open $V \subset Y$
the cohomology group $H^i(V, Rf_*K)$ is zero for $i < 0$ and is equal to
$H^0(V, H^0(Rf_*K))$ for $i = 0$. This proves (1).

\medskip\noindent
To prove (2) apply (1) to the complex $R\SheafHom(K, L)$ using
Lemma \ref{lemma-section-RHom-over-U} to do the translation.
\end{proof}

\begin{lemma}
\label{lemma-uniqueness}
In Situation \ref{situation-locally-given} assume
\begin{enumerate}
\item $X = \bigcup_{U \in \mathcal{B}} U$ and
for $U, V \in \mathcal{B}$ we have
$U \cap V = \bigcup_{W \in \mathcal{B}, W \subset U \cap V} W$,
\item for any $U \in \mathcal{B}$ we have $\Ext^i(K_U, K_U) = 0$
for $i < 0$.
\end{enumerate}
If a solution $(K, \rho_U)$ exists, then it is unique up to unique isomorphism
and moreover $\Ext^i(K, K) = 0$ for $i < 0$.
\end{lemma}

\begin{proof}
Let $(K, \rho_U)$ and $(K', \rho'_U)$ be a pair of solutions.
Let $f : X \to Y$ be the continuous map constructed
in Topology, Lemma \ref{topology-lemma-create-map-from-subcollection}.
Set $\mathcal{O}_Y = f_*\mathcal{O}_X$.
Then $K, K'$ and $\mathcal{B}$ are as in
Lemma \ref{lemma-vanishing-and-glueing} part (2).
Hence we obtain the vanishing of negative exts for $K$ and we see that
the rule
$$
V \longmapsto \Hom(K|_{f^{-1}V}, K'|_{f^{-1}V})
$$
is a sheaf on $Y$. As both $(K, \rho_U)$ and $(K', \rho'_U)$ are solutions
the maps
$$
(\rho'_U)^{-1} \circ \rho_U : K|_U \longrightarrow K'|_U
$$
over $U = f^{-1}(f(U))$ agree on overlaps. Hence we get a unique global
section of the sheaf above which defines the desired isomorphism
$K \to K'$ compatible with all structure available.
\end{proof}

\begin{remark}
\label{remark-uniqueness}
With notation and assumptions as in Lemma \ref{lemma-uniqueness}.
Suppose that $U, V \in \mathcal{B}$. Let $\mathcal{B}'$ be the set of
elements of $\mathcal{B}$ contained in $U \cap V$. Then
$$
(\{K_{U'}\}_{U' \in \mathcal{B}'},
\{\rho_{V'}^{U'}\}_{V' \subset U'\text{ with }U', V' \in \mathcal{B}'})
$$
is a system on the ringed space $U \cap V$
satisfying the assumptions of Lemma \ref{lemma-uniqueness}.
Moreover, both $(K_U|_{U \cap V}, \rho^U_{U'})$ and
$(K_V|_{U \cap V}, \rho^V_{U'})$ are solutions to this system.
By the lemma we find a unique isomorphism
$$
\rho_{U, V} : K_U|_{U \cap V} \longrightarrow K_V|_{U \cap V}
$$
such that for every $U' \subset U \cap V$, $U' \in \mathcal{B}$ the
diagram
$$
\xymatrix{
K_U|_{U'} \ar[rr]_{\rho_{U, V}|_{U'}} \ar[rd]_{\rho^U_{U'}} & &
K_V|_{U'} \ar[ld]^{\rho^V_{U'}} \\
& K_{U'}
}
$$
commutes. Pick a third element $W \in \mathcal{B}$. We obtain isomorphisms
$\rho_{U, W} : K_U|_{U \cap W} \to K_W|_{U \cap W}$ and
$\rho_{V, W} : K_U|_{V \cap W} \to K_W|_{V \cap W}$ satisfying
similar properties to those of $\rho_{U, V}$. Finally,
we have
$$
\rho_{U, W}|_{U \cap V \cap W} =
\rho_{V, W}|_{U \cap V \cap W} \circ \rho_{U, V}|_{U \cap V \cap W}
$$
This is true by the uniqueness in the lemma
because both sides of the equality are the unique isomorphism
compatible with the maps $\rho^U_{U''}$ and $\rho^W_{U''}$
for $U'' \subset U \cap V \cap W$, $U'' \in \mathcal{B}$.
Some minor details omitted.
The collection $(K_U, \rho_{U, V})$ is a descent datum
in the derived category for the open covering
$\mathcal{U} : X = \bigcup_{U \in \mathcal{B}} U$ of $X$.
In this language we are looking for ``effectiveness of the descent datum''
when we look for the existence of a solution.
\end{remark}

\begin{lemma}
\label{lemma-solution-in-finite-case}
In Situation \ref{situation-locally-given} assume
\begin{enumerate}
\item $X = U_1 \cup \ldots \cup U_n$ with $U_i \in \mathcal{B}$,
\item for $U, V \in \mathcal{B}$ we have
$U \cap V = \bigcup_{W \in \mathcal{B}, W \subset U \cap V} W$,
\item for any $U \in \mathcal{B}$ we have $\Ext^i(K_U, K_U) = 0$
for $i < 0$.
\end{enumerate}
Then a solution exists and is unique up to unique isomorphism.
\end{lemma}

\begin{proof}
Uniqueness was seen in Lemma \ref{lemma-uniqueness}. We may prove the lemma
by induction on $n$. The case $n = 1$ is immediate.

\medskip\noindent
The case $n = 2$.
Consider the isomorphism
$\rho_{U_1, U_2} : K_{U_1}|_{U_1 \cap U_2} \to K_{U_2}|_{U_1 \cap U_2}$
constructed in Remark \ref{remark-uniqueness}.
By Lemma \ref{lemma-glue} we obtain an object $K$ in $D(\mathcal{O}_X)$
and isomorphisms $\rho_{U_1} : K|_{U_1} \to K_{U_1}$ and
$\rho_{U_2} : K|_{U_2} \to K_{U_2}$ compatible with $\rho_{U_1, U_2}$.
Take $U \in \mathcal{B}$. We will construct an isomorphism
$\rho_U : K|_U \to K_U$ and we will leave it to the reader to verify
that $(K, \rho_U)$ is a solution. Consider the set $\mathcal{B}'$
of elements of $\mathcal{B}$ contained in either $U \cap U_1$ or contained in
$U \cap U_2$. Then $(K_U, \rho^U_{U'})$ is a solution for the system
$(\{K_{U'}\}_{U' \in \mathcal{B}'},
\{\rho_{V'}^{U'}\}_{V' \subset U'\text{ with }U', V' \in \mathcal{B}'})$
on the ringed space $U$.
We claim that $(K|_U, \tau_{U'})$ is another solution where
$\tau_{U'}$ for $U' \in \mathcal{B}'$ is chosen as follows:
if $U' \subset U_1$ then we take the composition
$$
K|_{U'} \xrightarrow{\rho_{U_1}|_{U'}}
K_{U_1}|_{U'} \xrightarrow{\rho^{U_1}_{U'}}
K_{U'}
$$
and if $U' \subset U_2$ then we take the composition
$$
K|_{U'} \xrightarrow{\rho_{U_2}|_{U'}}
K_{U_2}|_{U'} \xrightarrow{\rho^{U_2}_{U'}}
K_{U'}.
$$
To verify this is a solution use the property of the map $\rho_{U_1, U_2}$
described in Remark \ref{remark-uniqueness} and the compatibility of
$\rho_{U_1}$ and $\rho_{U_2}$ with $\rho_{U_1, U_2}$. Having said this
we apply Lemma \ref{lemma-uniqueness} to see that we obtain a unique
isomorphism $K|_{U'} \to K_{U'}$ compatible with the maps $\tau_{U'}$ and
$\rho^U_{U'}$ for $U' \in \mathcal{B}'$.

\medskip\noindent
The case $n > 2$. Consider the open subspace
$X' = U_1 \cup \ldots \cup U_{n - 1}$ and let $\mathcal{B}'$ be the set of
elements of $\mathcal{B}$ contained in $X'$. Then we find a system
$(\{K_U\}_{U \in \mathcal{B}'}, \{\rho_V^U\}_{U, V \in \mathcal{B}'})$
on the ringed space $X'$ to which we may apply our induction hypothesis.
We find a solution $(K_{X'}, \rho^{X'}_U)$.
Then we can consider the collection
$\mathcal{B}^* = \mathcal{B} \cup \{X'\}$ of opens of $X$ and we see that
we obtain a system
$(\{K_U\}_{U \in \mathcal{B}^*},
\{\rho_V^U\}_{V \subset U\text{ with }U, V \in \mathcal{B}^*})$.
Note that this new system also satisfies condition (3)
by Lemma \ref{lemma-uniqueness} applied to the solution $K_{X'}$.
For this system we have $X = X' \cup U_n$.
This reduces us to the case $n = 2$ we worked out above.
\end{proof}

\begin{lemma}
\label{lemma-glueing-increasing-union}
Let $X$ be a ringed space. Let $E$ be a well ordered set and let
$$
X = \bigcup\nolimits_{\alpha \in E} W_\alpha
$$
be an open covering with $W_\alpha \subset W_{\alpha + 1}$
and $W_\alpha = \bigcup_{\beta < \alpha} W_\beta$ if $\alpha$ is not
a successor. Let $K_\alpha$ be an object of $D(\mathcal{O}_{W_\alpha})$
with $\Ext^i(K_\alpha, K_\alpha) = 0$ for $i < 0$.
Assume given isomorphisms
$\rho_\beta^\alpha :  K_\alpha|_{W_\beta} \to K_\beta$ in
$D(\mathcal{O}_{W_\beta})$ for all $\beta < \alpha$ with
$\rho_\gamma^\alpha = \rho_\gamma^\beta \circ \rho^\alpha_\beta|_{W_\gamma}$
for $\gamma < \beta < \alpha$.
Then there exists an object
$K$ in $D(\mathcal{O}_X)$ and isomorphisms
$K|_{W_\alpha} \to K_\alpha$ for $\alpha \in E$
compatible with the isomorphisms $\rho_\beta^\alpha$.
\end{lemma}

\begin{proof}
In this proof $\alpha, \beta, \gamma, \ldots$ represent elements of $E$.
Choose a K-injective complex
$I_\alpha^\bullet$ on $W_\alpha$ representing $K_\alpha$.
For $\beta < \alpha$ denote $j_{\beta, \alpha} : W_\beta \to W_\alpha$
the inclusion morphism. Using transfinite recursion we will construct for all
$\beta < \alpha$ a map of complexes
$$
\tau_{\beta, \alpha} :
(j_{\beta, \alpha})_!I_\beta^\bullet
\longrightarrow
I_\alpha^\bullet
$$
representing the adjoint to the inverse of the isomorphism
$\rho^\alpha_\beta : K_\alpha|_{W_\beta} \to K_\beta$.
Moreover, we will do this in such that for
$\gamma < \beta < \alpha$ we have
$$
\tau_{\gamma, \alpha} = \tau_{\beta, \alpha} \circ
(j_{\beta, \alpha})_!\tau_{\gamma, \beta}
$$
as maps of complexes. Namely, suppose already given $\tau_{\gamma, \beta}$
composing correctly for all $\gamma < \beta < \alpha$.
If $\alpha = \alpha' + 1$ is a successor, then we choose any map of complexes
$$
(j_{\alpha', \alpha})_!I_{\alpha'}^\bullet \to I_\alpha^\bullet
$$
which is adjoint to the inverse of the isomorphism
$\rho^\alpha_{\alpha'} : K_\alpha|_{W_{\alpha'}} \to K_{\alpha'}$
(possible because $I_\alpha^\bullet$ is K-injective)
and for any $\beta < \alpha'$ we set
$$
\tau_{\beta, \alpha} = \tau_{\alpha', \alpha} \circ
(j_{\alpha', \alpha})_!\tau_{\beta, \alpha'}
$$
If $\alpha$ is not a successor, then
we can consider the complex on $W_\alpha$ given by
$$
C^\bullet = \colim_{\beta < \alpha} (j_{\beta, \alpha})_!I_\beta^\bullet
$$
(termwise colimit) where the transition maps of the sequence
are given by the maps $\tau_{\beta', \beta}$ for
$\beta' < \beta < \alpha$. We claim that $C^\bullet$
represents $K_\alpha$. Namely, for $\beta < \alpha$ the restriction
of the coprojection $(j_{\beta, \alpha})_!I_\beta^\bullet \to C^\bullet$
gives a map
$$
\sigma_\beta : I_\beta^\bullet \longrightarrow C^\bullet|_{W_\beta}
$$
which is a quasi-isomorphism: if $x \in W_\beta$ then looking
at stalks we get
$$
(C^\bullet)_x =
\colim_{\beta' < \alpha}
\left((j_{\beta', \alpha})_!I_{\beta'}^\bullet\right)_x =
\colim_{\beta \leq \beta' < \alpha} (I_{\beta'}^\bullet)_x
\longleftarrow
(I_\beta^\bullet)_x
$$
which is a quasi-isomorphism. Here we used that taking stalks
commutes with colimits, that filtered colimits are exact, and
that the maps $(I_\beta^\bullet)_x \to (I_{\beta'}^\bullet)_x$
are quasi-isomorphisms for $\beta \leq \beta' < \alpha$.
Hence $(C^\bullet, \sigma_\beta^{-1})$ is a solution to the
system $(\{K_\beta\}_{\beta < \alpha},
\{\rho^\beta_{\beta'}\}_{\beta' < \beta < \alpha})$.
Since $(K_\alpha, \rho^\alpha_\beta)$ is another solution
we obtain a unique isomorphism $\sigma : K_\alpha \to C^\bullet$
in $D(\mathcal{O}_{W_\alpha})$ compatible with all our maps, see
Lemma \ref{lemma-solution-in-finite-case}
(this is where we use the vanishing of negative ext groups).
Choose a morphism $\tau : C^\bullet \to I_\alpha^\bullet$
of complexes representing $\sigma$. Then we set
$$
\tau_{\beta, \alpha} = \tau|_{W_\beta} \circ \sigma_\beta
$$
to get the desired maps. Finally, we take $K$ to be the object of the derived
category represented by the complex
$$
K^\bullet = \colim_{\alpha \in E} (W_\alpha \to X)_!I_\alpha^\bullet
$$
where the transition maps are given by our carefully constructed
maps $\tau_{\beta, \alpha}$ for $\beta < \alpha$.
Arguing exactly as above we see that for all $\alpha$
the restriction of the coprojection determines an isomorphism
$$
K|_{W_\alpha} \longrightarrow K_\alpha
$$
compatible with the given maps $\rho^\alpha_\beta$.
\end{proof}

\begin{proof}
A pair $(K, \rho_U)$ is called a solution in the text above.
The uniqueness follows from Lemma \ref{lemma-uniqueness}.
If $X$ has a finite covering by elements of $\mathcal{B}$
(for example if $X$ is quasi-compact), then the theorem
is a consequence of Lemma \ref{lemma-solution-in-finite-case}.
In the general case we argue in exactly the same manner,
using transfinite induction and
Lemma \ref{lemma-glueing-increasing-union}.

\medskip\noindent
First we use transfinite recursion to choose opens $W_\alpha \subset X$
for any ordinal $\alpha$. Namely, we set $W_0 = \emptyset$.
If $\alpha = \beta + 1$ is a successor, then either $W_\beta = X$
and we set $W_\alpha = X$ or $W_\beta \not = X$ and we set
$W_\alpha = W_\beta \cup U_\alpha$ where
$U_\alpha \in \mathcal{B}$ is not contained in $W_\beta$.
If $\alpha$ is a limit ordinal we set
$W_\alpha = \bigcup_{\beta < \alpha} W_\beta$.
Then for large enough $\alpha$ we have $W_\alpha = X$.
Observe that for every $\alpha$ the open $W_\alpha$ is
a union of elements of $\mathcal{B}$. Hence if
$\mathcal{B}_\alpha = \{U \in \mathcal{B}, U \subset W_\alpha\}$, then
$$
S_\alpha = (\{K_U\}_{U \in \mathcal{B}_\alpha},
\{\rho_V^U\}_{V \subset U\text{ with }U, V \in \mathcal{B}_\alpha})
$$
is a system as in Lemma \ref{lemma-uniqueness} on the ringed space $W_\alpha$.

\medskip\noindent
We will show by transfinite induction that for every $\alpha$
the system $S_\alpha$ has a solution. This will prove the theorem
as this system is the system given in the theorem for large $\alpha$.

\medskip\noindent
The case where $\alpha = \beta + 1$ is a successor ordinal.
(This case was already treated in the proof of the lemma above
but for clarity we repeat the argument.)
Recall that $W_\alpha = W_\beta \cup U_\alpha$ for some
$U_\alpha \in \mathcal{B}$ in this case.
By induction hypothesis we have a solution
$(K_{W_\beta}, \{\rho^{W_\beta}_U\}_{U \in \mathcal{B}_\beta})$
for the system $S_\beta$.
Then we can consider the collection
$\mathcal{B}_\alpha^* = \mathcal{B}_\alpha \cup \{W_\beta\}$
of opens of $W_\alpha$ and we see that we obtain a system
$(\{K_U\}_{U \in \mathcal{B}_\alpha^*},
\{\rho_V^U\}_{V \subset U\text{ with }U, V \in \mathcal{B}_\alpha^*})$.
Note that this new system also satisfies condition (3)
by Lemma \ref{lemma-uniqueness} applied to the solution $K_{W_\beta}$.
For this system we have $W_\alpha = W_\beta \cup U_\alpha$.
This reduces us to the case handled in
Lemma \ref{lemma-solution-in-finite-case}.

\medskip\noindent
The case where $\alpha$ is a limit ordinal. Recall that
$W_\alpha = \bigcup_{\beta < \alpha} W_\beta$ in this case.
For $\beta < \alpha$ let
$(K_{W_\beta}, \{\rho^{W_\beta}_U\}_{U \in \mathcal{B}_\beta})$
be the solution for $S_\beta$.
For $\gamma < \beta < \alpha$ the restriction
$K_{W_\beta}|_{W_\gamma}$ endowed with the maps
$\rho^{W_\beta}_U$, $U \in \mathcal{B}_\gamma$
is a solution for $S_\gamma$. By uniqueness we get unique isomorphisms
$\rho_{W_\gamma}^{W_\beta} : K_{W_\beta}|_{W_\gamma} \to K_{W_\gamma}$
compatible with the maps $\rho^{W_\beta}_U$ and $\rho^{W_\gamma}_U$
for $U \in \mathcal{B}_\gamma$. These maps compose in the correct manner,
i.e., $\rho_{W_\delta}^{W_\gamma} \circ \rho_{W_\gamma}^{W_\beta}|_{W_\delta}
= \rho^{W_\delta}_{W_\beta}$ for $\delta < \gamma < \beta < \alpha$.
Thus we may apply Lemma \ref{lemma-glueing-increasing-union}
(note that the vanishing of negative exts is true for
$K_{W_\beta}$ by Lemma \ref{lemma-uniqueness} applied
to the solution $K_{W_\beta}$)
to obtain $K_{W_\alpha}$ and isomorphisms
$$
\rho_{W_\beta}^{W_\alpha} :
K_{W_\alpha}|_{W_\beta}
\longrightarrow
K_{W_\beta}
$$
compatible with the maps $\rho_{W_\gamma}^{W_\beta}$ for
$\gamma < \beta < \alpha$.

\medskip\noindent
To show that $K_{W_\alpha}$ is a solution we still need to construct the
isomorphisms $\rho_U^{W_\alpha} : K_{W_\alpha}|_U \to K_U$ for
$U \in \mathcal{B}_\alpha$ satisfying certain compatibilities.
We choose $\rho_U^{W_\alpha}$ to be the unique map such that
for any $\beta < \alpha$ and any $V \in \mathcal{B}_\beta$
with $V \subset U$ the diagram
$$
\xymatrix{
K_{W_\alpha}|_V \ar[r]_{\rho_U^{W_\alpha}|_V}
\ar[d]_{\rho_{W_\beta}^{W_\alpha}|_V}
& K_U|_V \ar[d]^{\rho_U^V} \\
K_{W_\beta} \ar[r]^{\rho_V^{W_\beta}}
& K_V
}
$$
commutes. This makes sense because
$$
(\{K_V\}_{V \subset U, V \in \mathcal{B}_\beta\text{ for some }\beta < \alpha},
\{\rho_V^{V'}\}_{V \subset V'\text{ with }V, V' \subset U
\text{ and }V, V' \in \mathcal{B}_\beta\text{ for some }\beta < \alpha})
$$
is a system as in Lemma \ref{lemma-uniqueness} on the ringed space $U$
and because $(K_U, \rho^U_V)$ and
$(K_{W_\alpha}|_U,  \rho_V^{W_\beta}\circ \rho_{W_\beta}^{W_\alpha}|_V)$
are both solutions for this system. This gives existence and uniqueness.
We omit the proof that these
maps satisfy the desired compatibilities (it is just bookkeeping).
\end{proof}

\end{document}
